\documentclass[10pt]{amsart}
\usepackage[a4paper,margin=22mm]{geometry}
\usepackage[no-math]{fontspec}
\setmainfont{Latin Modern Roman}[SmallCapsFont={lmromancaps10-regular.otf}]
\usepackage{amsmath,amssymb,amsthm,mathtools,hyperref,bbm,graphicx}
\usepackage[numbers,square]{natbib}
\newtheorem{theo}{Theorem}[section]
\newtheorem{lemm}[theo]{Lemma}
\newtheorem{prop}[theo]{Proposition}
\newtheorem{coro}[theo]{Corollary}
\theoremstyle{definition}
\newtheorem{defi}[theo]{Definition}
\newtheorem{exam}[theo]{Example}
\newtheorem{rema}[theo]{Remark}
\newtheorem{ques}[theo]{Question}
\newtheorem{conj}[theo]{Conjecture}
\newcommand{\longthanks}[1]{\section*{Acknowledgements}#1}
\newtheorem{lemma}[theo]{Lemma}
\newtheorem{question}[theo]{Question}
\newtheorem{conjecture}[theo]{Conjecture}
\newcommand{\ie}{i.e.}
\newcommand{\eg}{e.g.}
\emergencystretch=3em
\DeclareMathOperator\lk{lk}
\DeclareMathOperator\cost{cost}
\DeclareMathOperator{\Hilb}{Hilb}
\DeclareMathOperator{\Hom}{Hom}
\DeclareMathOperator{\Ext}{Ext}

\newcommand{\field}{\mathbbm{k}}
\newcommand{\NN}{{\mathbb N}}
\newcommand{\QQ}{{\mathbb Q}}
\newcommand{\CC}{{\mathbb C}}
\newcommand{\ZZ}{{\mathbb Z}}
\newcommand{\mideal}{\ensuremath{\mathfrak{m}}}
%\newcommand{\Hom}{\ensuremath{\mathrm{Hom}}\hspace{1pt}}
%\newcommand{\Ext}{\ensuremath{\mathrm{Ext}}\hspace{1pt}}

\graphicspath{{./figures/}}

\newcommand*{\mk}{\mkern -1mu}
\newcommand*{\Mk}{\mkern -2mu}
\newcommand*{\mK}{\mkern 1mu}
\newcommand*{\MK}{\mkern 2mu}

\newcommand*{\romanenumi}{\renewcommand*{\theenumi}{\roman{enumi}}}
\newcommand*{\Romanenumi}{\renewcommand*{\theenumi}{\Roman{enumi}}}
\newcommand*{\alphenumi}{\renewcommand*{\theenumi}{\alph{enumi}}}
\newcommand*{\Alphenumi}{\renewcommand*{\theenumi}{\Alph{enumi}}}

\title[Stanley--Reisner rings with a group action]{Stanley--Reisner rings of simplicial complexes with a free action by an abelian group}
\author{Connor Sawaske}
\date{Published 2018; DOI 10.5802/alco.29}
\begin{document}
\maketitle

\section{Introduction}

Since the 1970's, an extensive dictionary translating the topological and combinatorial properties of a simplicial complex $\Delta$ into algebraic properties of the associated Stanley--Reisner ring $\field[\Delta]$ has been constructed. For instance, the local cohomology modules $H_\mideal^i(\field[\Delta])$ of $\field[\Delta]$ have a beautiful interpretation due to Hochster (and later Gr\"abe) as topological invariants of $\Delta$ (see~\cite[Theorem~2]{Grabe}, \cite{Reisner}, or~\cite[Section~II.4]{St-96}). The following formulation of their results is the starting point of this paper (here $\cost_\Delta \sigma$ denotes the contrastar of a face $\sigma$ of $\Delta$, while $H^{i-1}(\Delta, \cost_\Delta \sigma)$ is the relative simplicial cohomology of the pair $(\Delta, \cost_\Delta \sigma)$ computed with coefficients in $\field$ and $H_\mideal^i(\field[\Delta])$ is the local cohomology module of $\field[\Delta]$):

\begin{theo}\label{Reisner}
Let $\Delta$ be a simplicial complex on vertex set $\{1, \ldots, n\}$ and let $\field$ be a field. Then
\[
H_\mideal^i(\field[\Delta])_{-j}\cong \bigoplus_{U\in T(\Delta)_j}H^{i-1}(\Delta, \cost_\Delta s(U))
\]
as vector spaces over $\field$, where $T(\Delta)_j=\{U\in \NN^n: s(U)\in \Delta$ and $|U|=j\}$, and for $U=(u_1, \ldots, u_n)\in\NN^n$, $s(U)=\{k:u_k>0\}$ and $|U|=\sum_{k=1}^nu_k$.
\end{theo}

On the other hand, the effect on the Stanley--Reisner ring of a simplicial complex being endowed with a group action has also been studied in some depth (see~\cite[Section~III.8]{St-96} and~\cite{AdinThesis}). Our primary goal is to bring these two topics together by studying the additional structure on $H_\mideal^i(\field[\Delta])$ that appears when $\Delta$ admits a group action. As it turns out, the topological invariants introduced by the group action dictate a more detailed description of $H_\mideal^i(\field[\Delta])$, providing our main result.

Our secondary goal is to apply this new decomposition of $H_\mideal^i(\field[\Delta])$ to extend two separate classes of previously-existing results describing the $h$-vectors (perhaps the most widely-recognized and studied combinatorial invariants) of certain simplicial complexes. One of these classes of results examines Buchsbaum simplicial complexes: using Theorem~\ref{Reisner}, Schenzel was able to calculate the Hilbert series of the quotient of $\field[\Delta]$ by a linear system of parameters in the case that $\Delta$ is Buchsbaum. The culmination of this line of study was the following theorem appearing in~\cite{Schenzel}, providing an algebraic interpretation of the $h$-vector of $\Delta$.

\begin{theo}[Schenzel]\label{Schenzel}
Let $\Delta$ be a $(d-1)$-dimensional Buchsbaum simplicial complex, and let $\Theta$ be a linear system of parameters for $\Delta$. Then
\[
\dim_\field \left(\field[\Delta]/\Theta\field[\Delta]\right)_i= h_i(\Delta)+\binom{d}{i}\sum_{j=0}^{i-1}(-1)^{i-j-1}\beta_{j-1}(\Delta),
\]
where $\beta_{j-1}(\Delta)$ denotes the reduced simplicial Betti number of $\Delta$ computed over $\field$.
\end{theo}

Since any dimension must be non-negative, the theorem provides lower bounds for the entries of $h(\Delta)$ in terms of some topological invariants of $\Delta$. More recently, these bounds have been further lowered by the use of socles (\cite{NS-Socle}) and the sigma module $\Sigma(\Theta; \field[\Delta])$ (originally introduced by Goto in~\cite{Goto}, and whose definition is provided in Section~\ref{sect:preliminaries}).

\begin{theo}[Murai--Novik--Yoshida]\label{NS-Socle}
Let $\Delta$ be a $(d-1)$-dimensional Buchsbaum simplicial complex
and let $\Theta$ be a linear system of parameters for $\Delta$. Then
\[
\dim_\field \left(\field[\Delta]/\Sigma(\Theta; \field[\Delta])\right)_i= h_i(\Delta)+\binom{d}{i}\sum_{j=0}^{i}(-1)^{i-j-1}\beta_{j-1}(\Delta).
\]
\end{theo}

Furthermore, the symmetry appearing in the $h$-vector described by Gr\"abe in~\cite{GrabeDS} was reformulated by Novik and Murai in~\cite[Proposition~1.1]{FaceNumbers}, and using Theorem~\ref{NS-Socle} it may be described as follows:

\begin{theo}[Murai--Novik--Yoshida]\label{Duality}
Let $\Delta$ be a $(d-1)$-dimensional %Buchsbaum simplicial complex 
orientable $\mathbbm{k}$-homology manifold
and let $\Theta$ be a linear system of parameters for $\Delta$. Then
\[
\dim_\field \left(\field[\Delta]/\Sigma(\Theta; \field[\Delta])\right)_i=\dim_\field \left(\field[\Delta]/\Sigma(\Theta; \field[\Delta])\right)_{d-i}
\]
for $i=0, \ldots, d$.
\end{theo}

The second class of results that we will enlarge deals with the $h$-vectors of complexes admitting specific group actions by $\ZZ/p\ZZ$ in the Cohen--Macaulay case. In particular, the following two theorems (\cite[Theorem~3.2]{St-87} and~\cite[Theorem~3.3]{AdinThesis}, respectively) provide impressive bounds on the $h$-vector of such a complex and are ripe for extensions to more general complexes:

\begin{theo}[Stanley]\label{StanleyCMVeryFree}
Let $\Delta$ be a $(d-1)$-dimensional Cohen--Macaulay (over $\CC$) simplicial complex admitting a very free action by $\ZZ/p\ZZ$ with $p$ a prime. Then $ h_i(\Delta)\ge \binom{d}{i}$ if $i$ is even and $h_i(\Delta)\ge (p-1)\binom{d}{i}$ if $i$ is odd.
\end{theo}

\begin{theo}[Adin]\label{AdinCMFree}
Let $\Delta$ be a $(d-1)$-dimensional Cohen--Macaulay (over $\CC$) simplicial complex with a free action of $\ZZ/p\ZZ$ with $p$ a prime such that $d$ is divisible by $p-1$. Then
\[
\sum_{i=0}^dh_i(\Delta)\lambda^i\ge (1+\lambda+\cdots +\lambda^{p-1})^{d/(p-1)},
\]
where the inequality holds coefficient-wise.
\end{theo}

In summary, our new results are the following (as much notation must be introduced before explicit statements may be provided, we present here a brief overview):
\begin{itemize}
%\itemsep0em
\item \emph{A new version of Hochster and Gr\"abe's theorem:} We examine how a group action leads to a special decomposition of $H_\mideal^i(\field[\Delta])$, and we demonstrate how the isomorphism of Theorem~\ref{Reisner} respects this decomposition (Theorem~\ref{ReisnerAnalog}). Later, for the $G=\ZZ/p\ZZ$ case, we encounter a piece of the local cohomology of $\field[\Delta]$ that corresponds to the singular cohomology of the quotient $|\Delta|/G$.
\item \emph{Applications to $h$-vectors:} The dimension calculations of Theorems~\ref{Schenzel} and~\ref{NS-Socle} are refined and specialized to the case of a Buchsbaum complex admitting a free action by a cyclic group of prime order (Theorems~\ref{SchenzelAnalog} and~\ref{SigmaModuleThm}). Using this refinement, Theorems~\ref{StanleyCMVeryFree} and~\ref{AdinCMFree} are generalized to the setting of Buchsbaum complexes; in some cases, the expressions are identical (Section~\ref{Inequalities}). Lastly, in Theorems~\ref{DSTheorem} and~\ref{DSTheorem2}, we exhibit a symmetry in the dimensions of finely graded pieces of some Artinian reductions of Stanley--Reisner rings of certain orientable homology manifolds without boundary admitting a free group action by a cyclic group of prime order in the flavor of Theorem~\ref{Duality}.
\end{itemize}
The paper is organized as follows. In Section~\ref{sect:preliminaries}, we introduce notation, review classical results, and discuss in some depth the theory of graded rings and modules. In Section~\ref{ReisnerSection} we study the local cohomology modules $H_\mideal^i(\field[\Delta])$ and provide our main new result. Section~\ref{Applications} is devoted to calculating a finely-graded Hilbert series of an Artinian reduction of the Stanley--Reisner ring $\field[\Delta]$, providing many applications of our main theorem. We close with comments and questions in Section~\ref{comments}.


\section{Preliminaries} \label{sect:preliminaries}

\subsection{Combinatorics and topology}\label{subsect:topo}

A \emph{simplicial complex} $\Delta$ on a finite vertex set $V$ is a collection of subsets of $V$ that is closed under inclusion. The elements of $\Delta$ are called \emph{faces}, and the maximal faces (with respect to inclusion) are \emph{facets}. The \emph{dimension} of a face $\sigma$ is defined by $\dim \sigma :=|\sigma|-1$, and the \emph{dimension} of $\Delta$ is defined by $\dim \Delta := \max\{\dim \sigma: \sigma\in \Delta\}$. The faces of dimension zero are called \emph{vertices}. If all facets of $\Delta$ have the same dimension, then we say that $\Delta$ is \emph{pure}.

Given a face $\sigma$ of $\Delta$, we define the \emph{contrastar} of $\sigma$ in $\Delta$ by
\[
\cost_\Delta \sigma := \{\tau\in \Delta: \sigma\not\subset\tau\}.
\]
Similarly, the \emph{link} of $\sigma$ in $\Delta$ is
\[
\lk_\Delta \sigma := \{\tau\in \Delta: \sigma\cup \tau \in \Delta, \sigma\cap \tau = \emptyset\}.
\]
We define the $f$-\emph{vector} $f(\Delta)$ of a $(d-1)$-dimensional simplicial complex $\Delta$ by
\[
f(\Delta)=(f_{-1}(\Delta), f_0(\Delta), \ldots, f_{d-1}(\Delta)),
\]
where $f_i(\Delta)$ is the number of $i$-dimensional faces of $\Delta$. The $h$-\emph{vector} $h(\Delta)$ is then defined by $h(\Delta)=(h_0(\Delta), h_1(\Delta), \ldots, h_d(\Delta))$, where
\[
h_i(\Delta)=\sum_{j=0}^i(-1)^{i-j}\binom{d-j}{i-j}f_{j-1}(\Delta).
\]


Let $\field$ be a field, and let $\tilde{H}^i(\Delta)$ be the $i$-th reduced simplicial cohomology group of $\Delta$ with coefficients in $\field$. If $\Gamma$ is a subcomplex of $\Delta$, then $H^i(\Delta, \Gamma)$ is the $i$-th relative cohomology group of the pair $(\Delta, \Gamma)$ with coefficients in $\field$. In the case that $\Gamma=\{\emptyset\}$, this is the same as the reduced cohomology group $\tilde{H}^i(\Delta)$. Denote the $i$\emph{-th (reduced) Betti number} of $\Delta$ over $\field$ by
\[
\beta_i(\Delta):=\dim_\field \widetilde{H}^i(\Delta).
\]

We call a complex $\Delta$ \emph{Cohen--Macaulay} if $\beta_i(\lk_\Delta \sigma)=0$ for all faces $\sigma$ and all $i<\dim \lk_\Delta \sigma$. Similarly, we call a complex \emph{Buchsbaum} if $\Delta$ is pure and $\beta_i(\lk_\Delta \sigma)=0$ for all faces $\sigma\not=\emptyset$ and all $i<\dim \lk_\Delta \sigma$. If $\Delta$ is $(d-1)$-dimensional and the link of each non-empty face $\sigma$ of $\Delta$ has the homology of a $(d-|\sigma|-1)$-sphere, then we say that $\Delta$ is a $\field$\emph{-homology manifold}; furthermore, if $\beta_{\dim\Delta}(\Delta)$ is equal to the number of connected components of $\Delta$, then we say that $\Delta$ is \emph{orientable}.

Now let $G$ be a finite group and suppose $G$ acts on a simplicial complex $\Delta$ (in particular, each $g\in G$ acts as a simplicial isomorphism on $\Delta$). We say that $G$ acts \emph{freely} if $g\cdot \sigma \not= \sigma$ for all faces $\sigma\not=\emptyset$ and all non-identity elements $g\in G$ (we will always take our actions to be defined on the left). If $G$ acts freely on $\Delta$ and, additionally, $(g\cdot\{v\})\cup\{v\}$ is not an edge of $\Delta$ for every vertex $v\in V$ and every non-identity element $g\in G$, then we say that $G$ acts \emph{very freely}.

In the case that $G=\ZZ/2\ZZ$, free and very free actions are equivalent and we call $\Delta$ \emph{centrally-symmetric} when such an action by $G$ on $\Delta$ exists. When extending a free group action by $G$ to the geometric realization $|\Delta|$ of $\Delta$, the action obtained is a covering space action (see~\cite[Section~1.3]{Hatcher}).

If $\Delta$ admits a group action by $G$, then this extends to an action on the reduced (simplicial) chain complex $\tilde{C}_\bullet(\Delta)$ by permuting its basis elements in accordance with the group action. This extends further to the cochain complex $\tilde{C}^\bullet(\Delta):=\Hom_\field(\tilde{C}_\bullet(\Delta), \field)$ in the usual way: if $\hat{\sigma}\in \tilde{C}^i(\Delta)$ is the basis element defined by $\hat{\sigma}(\tau)=0$ for $\sigma\not=\tau$ and $\hat{\sigma}(\sigma)=1$, then
\[
(g\cdot\hat{\sigma})(\tau)=\hat{\sigma}(g^{-1}\cdot\tau),
\]
and hence $g\cdot\hat{\sigma}=\widehat{(g\cdot\sigma)}$.

Lastly, given a vector $U=(u_1, \ldots, u_n)\in \NN^n$, let $s(U)=\{i:u_i>0\}$ be the support of $U$ and let $|U|=\sum_{i=1}^n u_j$ be the $L^1$-norm of $U$. If $\Delta$ is a simplicial complex on vertex set $[n]$, then denote
\[
T(\Delta)_j=\{U\in \NN^n:\text{$s(U)\in \Delta$ and $|U|=j$}\}.
\]

\subsection{Modules and group actions}\label{sect:mods}

Assume from now on that the field $\field$ is an extension of $\CC$. If $G$ is a finite abelian group and $M$ is a $\field[G]$-module, then $M$ admits a decomposition into isotypic components according to the action of $G$. That is,
\[
M=\bigoplus_{\chi\in\hat{G}} M^\chi
\]
where $\hat{G}$ is the group of irreducible characters $\chi$ of $G$ and
\[
M^\chi := \{ m\in M: g\cdot m = \chi(g) m\text{ for all $g\in G$}\}.
\]
Of vital importance will be the consideration of direct sums of simplicial cohomology modules of the form
\[
\bigoplus_{U\in T(\Delta)_j}H^{i-1}(\Delta, \cost_\Delta s(U)),
\]
which inherit a $G$-action from $\Delta$ as in the previous section. In particular, Theorem~\ref{ReisnerAnalog} will consider the isotypic components of modules of this form. In the $j=0$ case, we set
\begin{equation}\label{bettinotation}
\beta_i(\Delta)^\chi:=\dim_\field \widetilde{H}^i(\Delta)^\chi.
\end{equation}
These refined Betti numbers will be one of the main invariants considered in Sections~4 and 5 of this paper.


Now let $A$ be the polynomial ring $\field[x_1, \ldots, x_n]$, graded by degree, and let $M$ be a finitely-generated $\ZZ$-graded $A$-module, written as an abelian group by $M=\bigoplus_i M_i$. If $G$ acts on $M$ and each $M_i$ is $G$-invariant (\ie, $G\cdot M_i\subseteq M_i$), then each $M_i$ can itself be thought of as a $\field[G]$-module. As above,
\begin{equation}\label{decomp}
M_i=\bigoplus_{\chi\in\hat{G}} M_i^\chi
\end{equation}
where
\[
M_i^\chi=\{m\in M_i:g\cdot m = \chi(g)m\text{ for all $g\in G$}\}.
\]

Suppose $A$ is endowed with a degree-preserving action of $G$. Then the decomposition in~\eqref{decomp} makes $A$ into a $(\ZZ\times \hat{G})$-graded ring, since $a_1\in A_{i_1}^{\chi_1}$ and $a_2\in A_{i_2}^{\chi_2}$ implies that $a_1a_2\in A_{i_1+i_2}$ and that
\[
g\cdot(a_1a_2)=(g\cdot a_1)(g\cdot a_2)=\chi_1(g)a_1\chi_2(g)a_2=(\chi_1\chi_2)(g)a_1a_2
\]
for all $g\in G$. If, additionally,
\begin{equation}\label{compatible}
g\cdot(am)=(g\cdot a)(g\cdot m)
\end{equation}
for all $a\in A$, $g\in G$, and $m\in M$, then the decomposition in~\eqref{decomp} allows for a $(\ZZ\times \hat{G})$-grading for $M$, where $M_i^\chi$ is the component of $M$ consisting of all elements of degree $(i, \chi)$.

If $N$ is another $(\ZZ\times \hat{G})$-graded $A$-module satisfying the same conditions as $M$, then $G$ acts on $\Hom_A(M, N)$ by defining
\[
(g\cdot f)(m):=g\cdot f(g^{-1}\cdot m)
\]
for all $g\in G$, $f\in \Hom_A(M, N)$, and $m\in M$. If $m\in M_j^{\chi_1}$ and $f\in \Hom_A(M, N)_i^{\chi_2}$, then for all $g\in G$,
\[
\chi_2(g) f(m)=(g\cdot f)(m)=g\cdot f(g^{-1}\cdot m)=g\cdot f(\chi_1^{-1}(g)m)=\chi_1^{-1}(g)g\cdot f(m),
\]
so $g\cdot f(m)=(\chi_1\chi_2)(g)f(m)$ and $f(m)\in N_{i+j}^{\chi_1\chi_2}$. On the other hand, if $f\in \Hom_A(M, N)$ is such that $f(m)\in N_{i+j}^{\chi_1\chi_2}$ for all choices of $i$, $\chi_1$, and $m\in M_i^{\chi_1}$, then $g\cdot f=\chi_2(g) f$ for all $g\in G$. Hence,
\[
\Hom_A(M, N)_i^{\chi_1}=\{f\in\Hom_A(M, N): f(M_j^{\chi_2})\subseteq N_{i+j}^{\chi_1\chi_2}\text{ for all }(j,\chi_2)\in\ZZ\times \hat{G}\}.
\]
That is, the standard $(\ZZ\times \hat{G})$-grading on $\Hom_A(M, N)$ is consistent with the induced one above. In general, the set of graded $A$-module homomorphisms from $M$ to $N$ is a submodule of $\Hom_A(M, N)$. However, in our case of $M$ being finitely generated, the two submodules are equal (see~\cite[Section~II.11.6]{Bourbaki}). Unless stated otherwise, a $(\ZZ\times\hat{G})$-graded map of $A$-modules will refer to a homomorphism of degree~$0$.

Let $\mideal=(x_1, \ldots, x_n)$ denote the irrelevant ideal of $A$, let $H_\mideal^i(M)$ denote the $i$-th local cohomology module of $M$ with support in $\mideal$ (see~\cite{24Hours} for some basic properties of these modules), and denote $\mideal_j=(x_1^j, \ldots, x_n^j)$ (not to be confused with $M_j$, the $j$-th graded piece of $M$ as a $\ZZ$-graded $A$-module). We can trace the extra grading of $\Hom$ modules through to some standard local cohomology results using the following proposition.

\begin{prop}\label{extcohograding}
Let $G$ be a finite abelian group, and let $M$ be a $(\ZZ\times\hat{G})$-graded $A$-module with an action of $G$ satisfying~\eqref{compatible}. Then:
\begin{enumerate}%[label=\textup{(}\alph*\textup{)}]\alphenumi
\item \label{2.1.a} $\Ext_A^i(A/\mideal_j, M)$ and $H_\mideal^i(M)$ are $(\ZZ\times\hat{G})$-graded $A$-modules for all $i$ and all $j\ge 1$.
\item \label{2.1.b} The canonical maps $\psi_{i, j}^M:\Ext_A^i(A/\mideal_j, M)\to H_\mideal^i(M)$ are $(\ZZ\times\hat{G})$-graded.
\item \label{2.1.c} If $N$ satisfies the same conditions as $M$ and $f:M\to N$ is a map of\linebreak $(\ZZ\times\hat{G})$-graded $A$-modules, then the induced map $f^*:H_\mideal^i(M)\to H_\mideal^i(N)$ is also $(\ZZ\times\hat{G})$-graded.
\item \label{2.1.d} If $0\to L\to M\to N\to 0$ is a short exact sequence of $(\ZZ\times\hat{G})$-graded $A$-modules, then there is a long exact sequence
\[
\cdots\to H_\mideal^{i-1}(N)\xrightarrow{\delta} H_\mideal^i(L)\to H_\mideal^i(M)\to H_\mideal^i(N)
\xrightarrow{\delta} H_\mideal^{i+1}(L)\to\cdots
\]
of $(\ZZ\times\hat{G})$-graded $A$-modules.
\end{enumerate}
\end{prop}

\begin{rema}
While these statements are to be expected and follow the typical constructions, we could not find any all-encompassing reference for them. Since many of the modules and maps involved are vital to the proofs of Theorems~\ref{ReisnerAnalog} and~\ref{SchenzelAnalog}, the proofs have been included both for the sake of completeness and for a preview of what is to come.
\end{rema}

\begin{proof}
Let $K_\bullet^j$ denote the Koszul complex of $A$ with respect to the sequence $x_1^j, \ldots, x_n^j$. We view each term $K_t^j$ of $K_\bullet^j$ as
\[
K_t^j = \bigoplus_{1\le i_1<i_2<\cdots<i_t\le n}A(x_{i_1}^{j}\wedge x_{i_2}^{j}\wedge\cdots\wedge x_{i_t}^{j}),
\]
and we can naturally extend the action of $G$ on $A$ to an action of $G$ on $K_\bullet^j$ by defining
\[
g\cdot \left(f x_{i_1}^{j}\wedge\cdots\wedge x_{i_t}^{j}\right) = (g\cdot f)((g\cdot x_{i_1})^j\wedge\cdots\wedge (g\cdot x_{i_t})^j).
\]
Then equation~\eqref{compatible} holds for each $K_t^j$, making them $(\ZZ\times \hat{G})$-graded $A$-modules. The differential map on $K_\bullet^j$ preserves the $G$ action, so $\Hom_A(K_\bullet^j, M)$ is a cochain complex of $(\ZZ\times \hat{G})$-graded $A$-modules. Furthermore, the cohomology modules obtained from this cochain complex are also $(\ZZ\times \hat{G})$-graded $A$-modules (see~\cite[Proposition~II.11.3.3]{Bourbaki}).

Since $K_\bullet^j$ provides a projective resolution of $A/\mideal_j$,
\[
H^i(\Hom_A(K_\bullet^j, M))=\Ext_A^i(A/\mideal_j, M).
\]
The maps $\varphi_j :K_\bullet^{j+1}\to K_\bullet^j$ defined by $\varphi_j(x_{i_1}^{j+1}\wedge \cdots\wedge x_{i_t}^{j+1})=(x_{i_1}\!\cdots x_{i_t})x_{i_1}^{j}\wedge \cdots\wedge x_{i_t}^{j}$ also preserve the action of $G$, so the pullbacks
\[
\varphi_j^*:\Ext_A^i(A/\mideal_j, M)\to\Ext_A^i(A/\mideal_{j+1}, M)
\]
do as well. These maps provide a direct system of $(\ZZ\times \hat{G})$-graded $A$-modules in which
\[
H_\mideal^i(M):=\varinjlim\Ext_A^i(A/\mideal_j, M).
\]
By~\cite[Remark~II.11.3.3]{Bourbaki}, $H_\mideal^i(M)$ is a $(\ZZ\times \hat{G})$-graded $A$-module and the canonical maps $\psi_{i, j}^M:\Ext_A^i(A/\mideal_j, M)\to H_\mideal^i(M)$ also preserve the action of $G$. This proves~\eqref{2.1.a} and~\eqref{2.1.b}, from which~\eqref{2.1.c} follows almost immediately; if $m\in H_\mideal^i(M)$, choose $j$ and $m_j$ such that $m=\psi_{i, j}^M(m_j)$, and let
\[
f_j: \Ext_A^i(A/\mideal_j, M)\to \Ext_A^i(A/\mideal_j, N)
\]
be the map induced in $\Ext$ modules. It follows from the definitions that $f_j$ is a $(\ZZ\times\hat{G})$-graded map, while $f^*(m)=(f\circ \psi_{i, j}^M)(m_j)=(\psi_{i, j}^N\circ f_j)(m_j)$. Since all maps here are $(\ZZ\times\hat{G})$-graded, $f^*$ is as well.

For part~\eqref{2.1.d}, all that remains to be checked is that the connecting homomorphism $\delta$ is $(\ZZ\times\hat{G})$-graded. This follows from its standard construction, as it is defined as a composition of certain $(\ZZ\times\hat{G})$-graded maps.
\end{proof}

For any integer $a$, denote by $M[a]$ the shifted $\ZZ$-graded $A$-module defined by $M[a]_i=M_{i+a}$. If we are also given $\chi_1\in\hat{G}$, then we can define a new $G$-action $*$ on $M[a]$ by setting
\[
g*m=\chi_1^{-1}(g) (g\cdot m).
\]
for all $g\in G$. Denote by $M[a, \chi_1]$ the module $M[a]$ endowed with this new action of $G$. Then for all $g\in G$,
\[
\{m\in M_{i+a}:g\cdot m = (\chi_1\chi_2)(g)m\}=\{m\in M[a]_i:g * m = \chi_2(g) m\},
\]
hence $M[a, \chi_1]_i^{\chi_2}=M_{i+a}^{\chi_1\chi_2}$, and this makes $M[a, \chi_1]$ into a $(\ZZ\times \hat{G})$-graded $A$-module. Indeed, if $m\in M[a, \chi_1]_i^{\chi_2}$ and $f\in A_j^{\chi_3}$, then
\[
g*(fm)=\chi_1^{-1}(g)(g\cdot fm)=\chi_1^{-1}(g)(\chi_1\chi_2\chi_3)(g)fm=(\chi_2\chi_3)(g)fm,
\]
so that $fm\in M[a, \chi_1]_{i+j}^{\chi_2\chi_3}$. Moreover, if $f\in A_j^{\chi_2}$ then the multiplication map $\cdot f: M[-j, \chi_2^{-1}]\to M$ is a map of $(\ZZ\times G)$-graded $A$-modules (of degree~$0$), since if $m\in M[-j, \chi_2^{-1}]_i^{\chi_1}$, then
\[
f(g*m)=f\chi_2(g)(g\cdot m)=(g\cdot f)(g\cdot m)=g\cdot(fm).
\]

For the rest of this section, we will only be considering groups $G$ of the form $\ZZ/p\ZZ$ for a prime $p$. In this case, $G\cong\hat{G}$. Once we have fixed a generator $g$ for $G$ and a primitive $p$-th root of unity $\zeta$, there is a group isomorphism mapping $\chi\in \hat{G}$ to $j\in\ZZ/p\ZZ$ such that
\[
M^{\chi}=\{m\in M: g\cdot m = \zeta^jm\}.
\]
Hence, we write
\[
M=\bigoplus_{j=0}^{p-1}M^j
\]
where
\[
M^j:=\{m\in M: g\cdot m = \zeta^jm\}.
\]
With this identification, we view all of the modules constructed above as $(\ZZ\times G)$-graded modules, where $M_i^j$ is the component of $M$ in degree $(i, j)$. As in our definition~\eqref{bettinotation}, we set
\begin{equation}\label{bettiJnotation}
\beta_i(\Delta)^j:=\dim_\field \widetilde{H}^i(\Delta)^j.
\end{equation}

Thinking now of $M$ as a $(\ZZ\times G)$-graded vector space over $\field$, we can define the Hilbert series $\Hilb(M, \lambda, t)$ of $M$ by
\[
\Hilb(M, \lambda, t)=\sum_{(i, j)\in (\ZZ\times G)}(\dim_\field M_i^j)\lambda^it^j
\]
where $\lambda$ and $t$ are indeterminates with $t^p=1$.

If $M$ is of Krull dimension $d>0$, we call a system $\Theta = \theta_1, \theta_2, \ldots, \theta_d$ of homogeneous elements in $A$ a \emph{homogeneous system of parameters} (or an \text{h.s.o.p.}) for $M$ if $M/\Theta M$ is a finite-dimensional vector space over $\field$. If each $\theta_i\in A_1$, then we call $\Theta$ a \emph{linear system of parameters} (or an \text{l.s.o.p.}) for $M$. We say that $M$ is \emph{Cohen--Macaulay} if every l.s.o.p. is a regular sequence on $M$, and $M$ is \emph{Buchsbaum} if every l.s.o.p. satisfies
\[
(\theta_1, \ldots, \theta_{i-1})M:_M\theta_i = (\theta_1, \ldots, \theta_{i-1})M:_M\mideal
\]
for $i=1, \ldots, d$. Given any h.s.o.p.~$\Theta$ for $M$, we also have the notion of the sigma module $\Sigma(\Theta; M)$, defined by
\[
\Sigma(\Theta; M):=\Theta M + \sum_{i=1}^d\left((\theta_1, \ldots, \hat{\theta_i}, \ldots, \theta_d)M:_M\theta_i\right).
\]


\subsection{Stanley--Reisner rings}

Let $\Delta$ be a $(d-1)$-dimensional simplicial complex with vertex set $[n]:=\{1, \ldots, n\}$. Given $\sigma\subset [n]$, write $x_\sigma=\prod_{i\in \sigma}x_i$. The \emph{Stanley--Reisner ideal} of $\Delta$ is the ideal $I_\Delta$ of $A$ defined by
\[
I_\Delta=(x_\sigma:\sigma\subset [n], \sigma\not\in\Delta).
\]
The \emph{Stanley--Reisner ring} of $\Delta$ (over $\field$) is
\[
\field[\Delta]:= A/I_\Delta.
\]
We will always consider $\field[\Delta]$ as a module over $A$. The geometric notion of Buchsbaumness is tied algebraically to the Stanley--Reisner ring through the following vital result (found in~\cite{Schenzel}).

\begin{theo}[Schenzel]A pure simplicial complex $\Delta$ is Buchsbaum over $\field$ if and only if $\field[\Delta]$ is a Buchsbaum $A$-module.
\end{theo}

If $\Delta$ admits an action by the group $G=\ZZ/p\ZZ$, then this extends to an action on $\field[\Delta]$ and induces a $(\ZZ\times G)$-grading as detailed in the previous section. If this action is free, then for any $j$ and $i\ge 1$ we have
\[
\dim_\field \field[\Delta]_i^j=\frac{1}{p}\dim_\field \field[\Delta]_i.
\]
Also, $\dim_\field \field[\Delta]_0^0=1$ and $\dim_\field \field[\Delta]_0^j=0$ for $0<j<p$. As in~\cite[Section~3]{St-87}, this implies the following expression for the $(\ZZ\times G)$-graded Hilbert series of $\field[\Delta]$.

\begin{theo}\label{SRHilb}
Let $\Delta$ be a $(d-1)$-dimensional simplicial complex admitting a free action by the group $\ZZ/p\ZZ$. Then
\[
\Hilb(\field[\Delta], \lambda, t)=1+\frac{1}{p}\left[\frac{\sum_{i=1}^d h_i(\Delta)\lambda^i}{(1-\lambda)^d}-1\right](1+t+\cdots+t^{p-1}).
\]
\end{theo}

In fact, if $\Delta$ is centrally-symmetric (so $p=2$) and Cohen--Macaulay, then a certain Hilbert series provides a strong inequality bounding the $h$-vector of $\Delta$; see~\cite[Theorem~III.8.1]{St-96}. If $\Theta$ is an l.s.o.p. for $\Delta$, then we denote by $\field(\Delta; \Theta)$ the quotient $\field[\Delta]/\Theta\field[\Delta]$.

\begin{theo}[Stanley]
Let $\Delta$ be a $(d-1)$-dimensional centrally-symmetric Cohen--Macaulay simplicial complex, and suppose $\Theta=(\theta_1, \ldots, \theta_d)$ is an l.s.o.p. for $\field[\Delta]$ in which $\theta_i\in\field[\Delta]_1^1$ for $i=1, \ldots, d$. Then
\[
\Hilb(\field(\Delta; \Theta), \lambda, t)=\frac{1}{2}\left[(1-t)(1+\lambda)^d+(1+t)\sum_{i=0}^dh_i(\Delta)\lambda^i\right].
\]
\end{theo}


\section{The main theorem: a refinement of Hochster}\label{ReisnerSection}

For the rest of the paper, fix $\Delta$ to be a $(d-1)$-dimensional simplicial complex on vertex set $[n]$. Let $\field$ be an extension of $\CC$, and let $A$ be the polynomial ring $\field[x_1, \ldots, x_n]$. Recall that $T(\Delta)_j$ is the set $\{U\in \NN^n: s(U)\in \Delta$ and $ |U|=j\}$.

\begin{theo}\label{ReisnerAnalog}
Let $\Delta$ be a simplicial complex with a (not necessarily free) action by a finite abelian group $G$. Then the isomorphisms
\[
H_\mideal^i(\field[\Delta])_{-j}\cong \bigoplus_{U\in T(\Delta)_j}H^{i-1}(\Delta, \cost_\Delta s(U))
\]
of Theorem~\ref{Reisner} induce vector space isomorphisms
\[
H_\mideal^i(\field[\Delta])_{-j}^\chi\cong \left[\bigoplus_{U\in T(\Delta)_j}H^{i-1}(\Delta, \cost_\Delta s(U))\right]^{\chi}
\]
for each irreducible character $\chi$ of $G$.
\end{theo}

\begin{rema}
There are a multitude of proofs of Theorem~\ref{Reisner}, such as~\cite[Theorem~1]{Grabe} or~\cite[Corollary~4.4]{canonicalMap}. Our proof of the refinement above will be a modification of that which appears in~\cite{characterizations}.
\end{rema}

\begin{proof}
As in Section~\ref{sect:mods}, let $\mideal_{j+1}$ denote the ideal $(x_1^{j+1}, \ldots, x_n^{j+1})$ and let $K_\bullet^{j+1}$ denote the Koszul complex of $A$ with respect to the sequence $(x_1^{j+1}, \ldots, x_n^{j+1})$. If $C^\bullet(\Delta, \Gamma)$ denotes the relative simplicial cochain complex of the pair $(\Delta, \Gamma)$ with coefficients in $\field$, then we know from the proofs of Proposition~\ref{extcohograding} and~\cite[Theorem~1]{characterizations} that
\begin{align}\label{ExtAndCochainComplex}
H^i(\Hom_A(K_\bullet^{j+1}, \field[\Delta])_{-j})=\Ext_{A}^i(A/\mideal_{j+1}, \field[\Delta])_{-j}\cong \bigoplus_{\mathclap{{U\in T(\Delta)_j}}}H^{i-1}(C^\bullet(\Delta, \cost_\Delta s(U)))
\end{align}
as vector spaces over $\field$. Note that while~\cite[Theorem~1]{characterizations} is stated with respect to a $\ZZ^n$-grading, Reisner's proof in~\cite[Theorem~2]{Reisner} shows that the only non-acyclic components of the cochain complex $\Hom_A(K_\bullet^{j+1}, \field[\Delta])$ of $\ZZ$-graded $A$-modules also occur in the $\ZZ^n$-graded case.

Our first goal is to show that this isomorphism preserves isotypic components in the sense that it induces isomorphisms of the form
\begin{equation}\label{extiso}
\Ext_{A}^i(A/\mideal_{j+1}, \field[\Delta])_{-j}^\chi\cong \left[\bigoplus_{U\in T(\Delta)_j}H^{i-1}(C^\bullet(\Delta, \cost_\Delta s(U)))\right]^{\chi}
\end{equation}
for each irreducible character $\chi$ of $G$.



Given some set $\sigma=\{i_1, \ldots, i_t\}\subset [n]$ with $i_1<i_2<\cdots<i_t$, define $\overline{x}_\sigma:=x_{i_1}^{j+1}\wedge\cdots\wedge x_{i_t}^{j+1}$ in $K_t^{j+1}$ and recall that $x_\sigma := x_{i_1}\cdots x_{i_t}$ in $\field[\Delta]$. Likewise, if $U=(i_1, \ldots, i_n)\in \NN^n$, define $x_U:=x_{1}^{i_1}\cdots x_{n}^{i_n}$ in $\field[\Delta]$. If $\sigma\subset [n]$ is such that $|\sigma|=t$ and $U\in \NN^n$ is such that $s(U)\subset \sigma$ and $|U|=j$, define
\[
f_{\sigma}^U(\overline{x}_\tau)=
\begin{cases}
(x_\sigma)^{j+1}/x_U & \text{if } \tau=\sigma\\
0 & \text{if }\tau\not=\sigma.
\end{cases}
\]
Then $\Hom_A(K_t^{j+1}, \field[\Delta])_{-j}$ has a vector subspace with basis given by
\[
D^t=\{f_{\sigma}^U: |\sigma|=t, \text{ $\sigma\in \Delta$, $|U|=j$ and }s(U)\subset \sigma \}.
\]

For the right side of~\eqref{ExtAndCochainComplex}, given $\sigma\in \Delta$ with $|\sigma|=t$, as in Section~\ref{subsect:topo} let $\hat{\sigma}:C_{t-1}(\Delta)\to\field$ be the homomorphism defined by $\hat{\sigma}(\tau)=1$ if $\tau=\sigma$ and $\hat{\sigma}(\tau)=0$ otherwise and define
\[
\langle \sigma, U\rangle:=\hat{\sigma}+C^{t-1}(\cost_\Delta s(U)).
\]
Then a basis for
\[
\bigoplus_{U\in T(\Delta)_j}C^{t-1}(\Delta, \cost_\Delta s(U))
\]
is given by
\[
B=\{\langle \sigma, U\rangle : \sigma\in\Delta, \, |\sigma|=t, \, U\in T(\Delta)_j, \text{ and } s(U)\subseteq \sigma\},
\]
and Miyazaki's proof of Theorem~\ref{Reisner} shows that the direct sum of maps
\[
\psi:\bigoplus_{U\in T(\Delta)_j}C^{t-1}(\Delta, \cost_\Delta s(U))\to \Hom_A(K_t^{j+1}, \field[\Delta])_{-j}
\]
defined componentwise by
\[
\psi(\langle \sigma, U\rangle)=f_{\sigma}^U
\]
is the chain map inducing the isomorphism in cohomology in~\eqref{ExtAndCochainComplex}.

Now consider how $G$ acts on each of these cochain complexes. First, $g\cdot f_\sigma^U=f_{g\cdot \sigma}^{g\cdot U}$ from the definition of the action on $\Hom_A(K_\bullet^{j+1}, \field[\Delta])$. On the other hand,
\[
g\cdot\langle \sigma, U\rangle=\langle g \cdot \sigma, g\cdot U\rangle,
\]
where $G$ acts on $U\in\NN^n$ by permuting indices in accordance with the action of $G$ on the vertex set $[n]$. Hence, $\psi$ is a $G$-equivariant isomorphism and~\eqref{extiso} is established.

Finally, again referring to the proof of~\cite[Theorem~1]{characterizations}, the canonical maps
\[
\Ext_A^i(A/\mideal_{j+1}, \field[\Delta])\to H_\mideal^i(\field[\Delta])
\]
induce isomorphisms in $\ZZ$-graded components of degree at least $-j$. Since these are also maps of $(\ZZ\times \hat{G})$-graded $A$-modules by Proposition~\ref{extcohograding}, we obtain the desired isomorphism in the statement of the theorem.
\end{proof}


\section{Applications}\label{Applications}

\subsection{Hilbert series of Artinian reductions and an analog of Schenzel's formula}\label{SchenzelSection}
From now on, $\Delta$ is fixed to be Buchsbaum and $G=\ZZ/p\ZZ$ for some prime $p$ with a fixed generator $g$. We consider $A$ to be $(\ZZ\times G)$-graded and $\field[\Delta]$ to be a $(\ZZ\times G)$-graded $A$-module as in Section~\ref{sect:preliminaries}. In a way similar to~\cite[Section~8]{St-96}, if the action of $G$ is very free and $0\le \delta_i\le p-1$ for $i=1, \ldots, d$, then we can easily construct a l.s.o.p. $\Theta=\theta_1, \ldots, \theta_d$ for $\field[\Delta]$ in which $\theta_i\in A_1^{\delta_i}$ for $1\le i \le d$ as follows.

First, choose one face from each $G$-orbit of $\Delta$ and collect them into the set $\Delta_G$. Let $W$ be the set of vertices of $\Delta$ that are in $\Delta_G$, and choose functions $t_1, \ldots, t_d:W\to\field$ such that their restrictions to any subset of $W$ of size $d$ are linearly independent. Now extend $t_i$ to all of $[n]$ by setting
\[
t_i(g^k\cdot v)=\zeta^{-k\delta_i}t_i(v)
\]
for all $i$ and $k$, and let
\[
\theta_i=\sum_{v\in [n]}t_i(v)x_v.
\]
Then
\begin{align*}
g\cdot\theta_i &= \sum_{v\in [n]}t_i(v)x_{g\cdot v}=\sum_{v\in [n]}t_i(g^{-1}\cdot v)x_v=\sum_{v\in [n]}\zeta^{\delta_i}t_i(v)x_v,
\end{align*}
so $\theta_i\in A_1^{\delta_i}$ for $i=1, \ldots, d$. Furthermore, since no facet of $\Delta$ contains $\{g^i\cdot v, g^j\cdot v\}$ for any $v$ with $i\not\equiv j \mod p$, the system $\Theta=(\theta_1, \ldots, \theta_d)$ forms an l.s.o.p. that is homogeneous with respect to the $(\ZZ\times G)$-grading by~\cite[Lemma~III.2.4$\MK$(a)]{St-96}.

\begin{rema}\label{freeactionlsop}
The construction above is included purely for the sake of concreteness in the case of a very free action. In fact, as Adin shows in his thesis (\cite{AdinThesis}), it is possible to construct an l.s.o.p. with the prescribed properties above even in the case of a free action. The construction involves changing the field $\field$ to a particular extension of $\CC$. However, as we are now turning our focus to certain Hilbert series (which are invariant under field extensions), all of the results that follow hold for any extension of $\CC$ and are stated with this understanding in mind.
\end{rema}


Now that the existence of l.s.o.p.'s of the form above has been established, we may use them to prove the following theorem (recall the refined Betti number notation of definition~\eqref{bettiJnotation} and that $\field(\Delta; \Theta)=\field[\Delta]/\Theta\field[\Delta]$ for an l.s.o.p. $\Theta$).

\begin{theo}\label{SchenzelAnalog}
Let $\Delta$ be a $(d-1)$-dimensional Buchsbaum simplicial complex admitting a free group action by $G=\ZZ/p\ZZ$, let $0\le m\le p-1$ be some fixed degree, and let $\Theta$ be a $G$-homogeneous l.s.o.p. for $\field[\Delta]$ such that $\theta_i\in A_1^m$ for all $i$. Then the $(\ZZ\times G)$-graded Hilbert series of $\field(\Delta;\Theta)$ is given by
\begin{align*}
\Hilb(\field(\Delta; \Theta), \lambda, t)&=\sum_{i=0}^d\left[(-1)^i\binom{d}{i}t^{mi}+\left(\frac{1}{p}\sum_{k=0}^{p-1}t^k\right)\left(h_i(\Delta)+(-1)^{i+1}\binom{d}{i}\right)\right]\lambda^i\\
&\quad\,+\sum_{i=0}^d\binom{d}{i}\lambda^i\sum_{j=0}^{i-1}(-1)^{i-j-1}\left(\sum_{k=0}^{p-1}t^k\beta_{j-1}(\Delta)^{k-im} \right).
\end{align*}
\end{theo}

\begin{rema}
Recall that $t^p=1$ in this Hilbert series, and note that by setting $t=1$, we can recover Schenzel's Theorem~\ref{Schenzel}. Furthermore, we conclude that $h_i(\Delta) \equiv (-1)^i\binom{d}{i} \pmod p$ for all $i$.
\end{rema}

\begin{proof}
Suppose that $\Theta$ is an arbitrary l.s.o.p.\@ with $\theta_i\in A_1^{\delta_i}$ for $i=1, \ldots, d$. For $s=1, \ldots, d$, denote
\[
\field_s[\Delta]:=\field[\Delta]/(\theta_1, \ldots, \theta_s)\field[\Delta].
\]
Then (recall the shifted module discussion from Section~\ref{sect:mods}) we have exact sequences of $(\ZZ\times G)$-graded $A$-modules with degree-preserving maps of the form
\[
0\to Q(s)[-1, -\delta_s]\to \field_{s-1}[\Delta][-1, -\delta_s]\xrightarrow{\cdot\theta_s} \field_{s-1}[\Delta]\to \field_s[\Delta]\to 0,
\]
where
\[
Q(s)=0:_{\field_{s-1}[\Delta]}\theta_s=0:_{\field_{s-1}[\Delta]}\mideal = H_\mideal^0(\field_{s-1}[\Delta]).
\]
The second equality above follows from the definition of Buchsbaumness, and the third follows from the proof of~\cite[Proposition~I.2.1]{StVo}. On the level of Hilbert series, a standard argument yields the following equation:
\begin{multline*}
\Hilb(\field(\Delta; \Theta), \lambda, t)=\Hilb(\field[\Delta], \lambda, t)\prod_{i=1}^d(1-\lambda t^{\delta_i})\\+\sum_{s=1}^d\lambda t^{\delta_i}\Hilb(H_\mideal^0(\field_{s-1}[\Delta]), \lambda, t)\prod_{j=s+1}^d(1-\lambda t^{\delta_j}).
\end{multline*}
In the case $\theta_s\in A_1^m$ for $s=1, \ldots, d$, the equation above simplifies to
\begin{multline*}
\Hilb(\field(\Delta; \Theta), \lambda, t)=(1-\lambda t^m)^d\Hilb(\field[\Delta], \lambda, t)\\+\sum_{s=1}^d\lambda t^m(1-\lambda t^m)^{d-s}\Hilb(H_\mideal^0(\field_{s-1}[\Delta]), \lambda, t).
\end{multline*}
Analyzing the first term yields the following, using Theorem~\ref{SRHilb}:
\begin{align*}
(1-\lambda t^m)^d&\Hilb(\field[\Delta], \lambda, t)\\
&=(1-\lambda t^m)^d\left[1+\frac{1}{p}\left(\frac{\sum_{i=0}^d h_i(\Delta)\lambda^i}{(1-\lambda)^d}-1\right)(1+t+\cdots+t^{p-1})\right]\\
&=(1-\lambda t^m)^d+\frac{(1+t+\cdots+t^{p-1})}{p}\left[\sum_{i=0}^dh_i(\Delta)\lambda^i - (1-\lambda t^m)^d\right]\\
&=\sum_{i=0}^d\left[(-1)^i\binom{d}{i}t^{mi}+\frac{1}{p}\sum_{k=0}^{p-1}t^k\left(h_i(\Delta)+(-1)^{i+1}\binom{d}{i}\right)\right]\lambda^i.
\end{align*}
For the second term,
\[
H_\mideal^0(\field_{s-1}[\Delta])\cong\bigoplus_{i=0}^{s-1}\left(\bigoplus_{\binom{s-1}{i}}H_\mideal^i(\field[\Delta])[-i, -im]\right)
\]
by~\cite[Proposition~II.4.14$'$]{StVo} (though the stated result is for $\ZZ$-graded $A$-modules, the exact same proof works in the $(\ZZ\times G)$-graded case using Proposition~\ref{extcohograding}). By our Theorem~\ref{ReisnerAnalog}, $H^i_\mideal(\field[\Delta])[-i, -im]$ is concentrated in $\ZZ$-degree $i$ and has dimension $\beta_{i-1}(\Delta)^{k-im}$ in degree $(i, k)$. Hence,
\[
\Hilb(H_\mideal^0(\field_{s-1}[\Delta]), \lambda, t)=\sum_{i=0}^{s-1}\binom{s-1}{i}\lambda^i\left(\sum_{k=0}^{p-1}t^k\beta_{i-1}(\Delta)^{k-im}\right),
\]
and so the $\lambda^it^k$ coefficient of $\lambda t^m\Hilb(H_\mideal^0(\field_{s-1}[\Delta]), \lambda, t)$ is $\binom{s-1}{i-1}\beta_{i-2}(\Delta)^{k-im}$. Then the $\lambda^it^k$ coefficient of $(1-\lambda t^m)^{d-s}\lambda t^m\Hilb(H_\mideal^0(\field_{s-1}[\Delta]), \lambda, t)$ is given by
\[
\sum_{r=0}^i(-1)^r\binom{d-s}{r}\binom{s-1}{i-r-1}\beta_{i-r-2}(\Delta)^{k-im}.
\]
Now we sum over all values of $s$, setting $f(s)=\binom{d-s}{r}\binom{s-1}{i-r-1}$:
\begin{align*}
\sum_{s=1}^d\sum_{r=0}^i(-1)^rf(s)\beta_{i-r-2}(\Delta)^{k-im}&=\sum_{r=0}^i\left[(-1)^r\beta_{i-r-2}(\Delta)^{k-im}\sum_{s=1}^df(s)\right]\\
&=\binom{d}{i}\sum_{r=0}^i(-1)^r\beta_{i-r-2}(\Delta)^{k-im}\\
&=\binom{d}{i}\sum_{j=0}^{i-1}(-1)^{i-j-1}\beta_{j-1}(\Delta)^{k-im}.
\end{align*}
Here the second equality follows from the identity $\sum_{s=1}^{d}\binom{d-s}{r}\binom{s-1}{i-r-1}=\binom{d}{i}$ and the third follows from setting $j=i-r-1$.
\end{proof}

Of particular interest is the $G=\ZZ/2\ZZ$ case, in which many simplifications can be made to the expression in Theorem~\ref{SchenzelAnalog}. This results in the following corollary.

\begin{coro}\label{CsSchenzel}
Let $\Delta$ be a centrally-symmetric Buchsbaum complex of dimension $d-1$ with $\Theta=(\theta_1, \ldots, \theta_d)$ an l.s.o.p. for $\Delta$ such that $\theta_i\in A_1^m$ for all $i$ and some fixed $m$. Then
\[
\dim_\field \field(\Delta; \Theta)_i^k = \frac{1}{2}\left(h_i(\Delta)+(-1)^{i+k+mi}\binom{d}{i}\right)+\binom{d}{i}\sum_{j=0}^{i-1}(-1)^{i-j-1}\beta_{j-1}(\Delta)^{k-im}.
\]
\end{coro}

\subsection{The sigma module}

When only considering a $\ZZ$-grading on $\field[\Delta]$, it is possible to mod $\field(\Delta; \Theta)$ out by an additional submodule in order to get an even tighter bound on the $h$-vector of $\Delta$. In particular, the sigma module can be used to great effect as follows (\cite[Theorem~1.2]{BBMD}).

\begin{theo}[Murai--Novik--Yoshida]\label{BBMD}
Let $\Delta$ be a Buchsbaum simplicial complex of dimension $d-1$ and let $\Theta$ be an l.s.o.p. for $\Delta$. Then
\[
\dim_\field \left(\field[\Delta]/\Sigma(\Theta; \field[\Delta])\right)_i = h_i(\Delta)+\binom{d}{i}\sum_{j=0}^i(-1)^{i-j-1}\beta_{j-1}(\Delta).
\]
\end{theo}

Of course, we would like for analogous statements to hold for the $(\ZZ\times G)$-graded Hilbert series of $\field[\Delta]/\Sigma(\Theta; \field[\Delta])$. In order to even consider this series, we must first verify that this module is in fact $(\ZZ\times G)$-graded by establishing that the sigma module is $G$-invariant. This is accomplished by the following lemma, whose proof is nearly immediate and has been omitted.

\begin{lemma}
Let $\Delta$ be a Buchsbaum simplicial complex with a free action by $G$. If $\Theta=(\theta_1, \ldots, \theta_d)$ is an l.s.o.p. for $\Delta$ and each $\theta_i$ is homogeneous with respect to the $(\ZZ\times G)$-grading of $A$, then the sigma module $\Sigma (\Theta; \field[\Delta])$ is $G$-invariant.
\end{lemma}

Fortunately, the proof of our needed aspects of~\cite[Theorem~2.3]{BBMD} goes through essentially verbatim in the $(\ZZ\times G)$-graded case using Proposition~\ref{extcohograding}. One aspect not covered directly by the proposition (though it is considered in the proof) is that images and kernels of maps of $(\ZZ\times\hat{G})$-graded $A$-modules are themselves $(\ZZ\times\hat{G})$-graded; this follows from~\cite[Proposition~II.11.3.3]{Bourbaki}. Hence, we obtain the following proposition.

\begin{prop}\label{sigmaModuleProp}
Let $\Delta$ be a $(d-1)$-dimensional Buchsbaum simplicial complex admitting a free action by $G$ and let $\Theta=(\theta_1, \ldots, \theta_d)$ be an l.s.o.p. for $\Delta$ with $\theta_i\in A_1^m$ for all $i$. Then
\[
\Sigma(\Theta;\field[\Delta])/\Theta\field[\Delta]\cong\bigoplus_{i=0}^{d-1}\left(\bigoplus_{\binom{d}{i}}H_\mideal^{i}(\field[\Delta])[-i, -im]\right)
\]
as $(\ZZ\times G)$-graded $A$-modules. In particular,
\[
\dim_{\field}\left(\Sigma(\Theta; \field[\Delta])/\Theta\field[\Delta]\right)_i^k=\binom{d}{i}\beta_{i-1}(\Delta)^{k-im}.
\]
\end{prop}

This immediately establishes the following extension of Theorem~\ref{SchenzelAnalog}:

\begin{theo}\label{SigmaModuleThm}
Let $\Delta$ be a $(d-1)$-dimensional Buchsbaum simplicial complex admitting a free action by $G$ and let $\Theta$ be a $G$-homogeneous l.s.o.p. for $\field[\Delta]$ such that $\theta_i\in A_1^m$ for all $i$. Then
\begin{align*}
\Hilb(\field[\Delta]&/\Sigma(\Theta;\field[\Delta]), \lambda, t)=\\
&\sum_{i=0}^d\left[(-1)^i\binom{d}{i}t^{mi}+\left(\frac{1}{p}\sum_{k=0}^{p-1}t^k\right)\left(h_i(\Delta)+(-1)^{i+1}\binom{d}{i}\right)\right]\lambda^i\\
&+\sum_{i=0}^d\binom{d}{i}\lambda^i\sum_{j=0}^{i}(-1)^{i-j-1}\left(\sum_{k=0}^{p-1}t^k\beta_{j-1}(\Delta)^{k-im} \right).
\end{align*}
\end{theo}



\subsection{Inequalities}\label{Inequalities}
Theorem~\ref{BBMD} implies the following inequality bounding the $h$-vector of any Buchsbaum simplicial complex $\Delta$:
\begin{equation}\label{buchsbaumInequality}
h_i(\Delta)\ge \binom{d}{i}\sum_{j=0}^i(-1)^{i-j}\beta_{j-1}(\Delta).
\end{equation}
For Cohen--Macaulay complexes admitting a very free action by $G$, bounds on the $h$-vector may be obtained through Theorem~\ref{StanleyCMVeryFree}. In this section we present some similar inequalities by exploiting the $(\ZZ\times G)$-graded Hilbert series of Buchsbaum complexes admitting free group actions by $G$. To begin, setting $m=0$ and examining the $\lambda^it^{k}$ coefficients in Theorem~\ref{SigmaModuleThm} immediately allows us to bound the $h$-vector of $\Delta$ by
\begin{equation}\label{zeropart}
h_i(\Delta)\ge(p-1)(-1)^{i+1}\binom{d}{i}+p\binom{d}{i}\sum_{j=0}^i(-1)^{i-j}\beta_{j-1}(\Delta)^0
\end{equation}
and
\begin{equation}\label{nonzeropart}
h_i(\Delta)\ge(-1)^{i}\binom{d}{i}+p\binom{d}{i}\sum_{j=0}^i(-1)^{i-j}\beta_{j-1}(\Delta)^k
\end{equation}
for $k\not\equiv 0 \mod p$.

The inequalities above provide an immediate ``permutable'' version of~\eqref{buchsbaumInequality} that is worth recognizing.

\begin{coro}\label{multiset}
Let $\Delta$ be a $(d-1)$-dimensional Buchsbaum simplicial complex admitting a free action by the group $\ZZ/p\ZZ$. If $\mathcal{M}$ is a multiset of size $p-1$ on $\{1, \ldots, p-1\}$, then
\[
h_i(\Delta)\ge\binom{d}{i}\sum_{j=0}^i(-1)^{i-j}\beta_{j-1}(\Delta)^0+\binom{d}{i}\sum_{m\in\mathcal{M}}\sum_{j=0}^i(-1)^{i-j}\beta_{j-1}(\Delta)^m.
\]
\end{coro}

\begin{proof}
For each $m\in \mathcal{M}$, consider the corresponding inequality~\eqref{nonzeropart}. Add all of the inequalities together along with~\eqref{zeropart}, then divide by $p$.
\end{proof}

Note that by taking $\mathcal{M}=[p-1]$, we obtain the original bound~\eqref{buchsbaumInequality}. On the other hand, by considering the degrees of the group representations produced by $G$ acting on $H^i(\Delta)$, we can also obtain some extensions of Theorem~\ref{StanleyCMVeryFree}.

\begin{coro}
Let $\Delta$ be a $(d-1)$-dimensional Buchsbaum simplicial complex admitting a free action by $G$. Then:
\begin{enumerate}%[label=\textup{(}\alph*\textup{)}]\alphenumi
\item \label{4.9.a} If $\beta_i(\Delta)=0$ for $i\le j$, then $h_i(\Delta)\ge\binom{d}{i}$ for $i\le j+2$.
\item \label{4.9.b} If $\beta_i(\Delta)<(p-1)$ for $i\le j$, then $h_i(\Delta)\ge (-1)^i\binom{d}{i}$ for $i\le j+2$.
\end{enumerate}
\end{coro}

\begin{proof}
The statement in~\eqref{4.9.a} is an immediate consequence of the inequalities in~\eqref{zeropart} and~\eqref{nonzeropart}. For~\eqref{4.9.b}, consider the reduced cohomology group $\tilde{H}^i(\Delta; \QQ)$ with coefficients in $\QQ$, which satisfies
\[
\beta_i(\Delta)=\dim_\QQ \tilde{H}^i(\Delta; \QQ).
\]
Furthermore,
\[
\tilde{H}^i(\Delta; \field)\cong \tilde{H}^i(\Delta; \QQ)\otimes_\QQ \field
\]
as representations of $G$ by taking $\field$ to have a trivial action.

There are only two irreducible representations of $G$ over $\QQ$: the $1$-dimensional trivial representation and a $(p-1)$-dimensional non-trivial representation. Hence, if $\dim_\QQ \tilde{H}^i(\Delta; \QQ)<p-1$ then $G$ acts trivially on $\tilde{H}^i(\Delta; \QQ)$. From the above isomorphism, $G$ acts trivially on $\tilde{H}^i(\Delta;\field)$ as well. That is, $\beta_i(\Delta)^r =0$ for $1\le r\le p-1$ and $i\le j$, and the result follows from inequality~\eqref{nonzeropart}.
\end{proof}

It is also worth examining the properties of the quotient $\Delta/G$. This quotient has elements corresponding to orbits of the faces of $\Delta$ under the action of $G$, and it can always be made into a poset with zero element $\{\emptyset\}$ (assuming $\Delta\not=\emptyset$) under the ordering defined by setting $x\preceq y$ if $\sigma\subseteq \tau$ for some $\sigma\in x$ and $\tau\in y$. This poset is ranked by setting the rank of an orbit to be the cardinality of any of its members.

In some cases, $\Delta/G$ is itself isomorphic to the face poset of a simplicial complex (for instance, when $\ZZ/2\ZZ$ acts by rotation on the boundary of a hexagon, the quotient is isomorphic to the face poset of the boundary of a triangle). When this happens, we naturally view $\Delta/G$ as the corresponding simplicial complex. Less optimistically, $\Delta/G$ may have the property that every interval $[\{\emptyset\}, x]$ is isomorphic to a boolean lattice. In this case, we say that $\Delta/G$ is a simplicial poset. Many invariants of simplicial complexes have natural extensions to simplicial posets (see~\cite{StanleyPosets} for a nice overview). In particular, the $h$-vector of a simplicial poset is a well-studied object that our next proposition examines.

\begin{prop}
Let $\Delta$ be a Buchsbaum simplicial complex admitting a free group action by $G$. If $\Delta/G$ is a simplicial poset, then
\[
h_i(\Delta/G)\ge (-1)^i\binom{d}{i}+\binom{d}{i}\sum_{j=0}^i(-1)^{i-j}\beta_{j-1}(\Delta)^k
\]
for $k\not\equiv 0 \mod p$.
\end{prop}

\begin{proof}
A straightforward calculation shows that
\begin{equation}\label{hVectorOfQuotient}
h_i(\Delta)=(p-1)(-1)^{i+1}\binom{d}{i}+p h_i(\Delta/G).
\end{equation}
Now combine~\eqref{nonzeropart} with~\eqref{hVectorOfQuotient} and divide by $p$.
\end{proof}

As examples, $\Delta/G$ is always a simplicial poset when $\Delta$ is the order complex of a poset (\cite[Section~6]{Garsia}), or, more generally, when $\Delta$ is balanced and the action of $G$ is color-preserving (\cite{Reiner}). Note that this proposition provides a ``Buchsbaum-like'' bound on the $h$-vector of $\Delta/G$ in the style of equation~\eqref{nonzeropart} without any direct consideration of whether or not $\Delta/G$ is itself Buchsbaum.



\subsection{Dehn--Sommerville relations} Here we present a new results akin to Theorem~1.4, relating the symmetric values of the dimensions $\dim_\field \left(\field[\Delta]/\Sigma(\Theta;\field[\Delta])\right)_i^j$. First, let $\tilde{\chi}(\Delta)$ denote the reduced Euler characteristic of a simplicial complex $\Delta$, defined by
\[
\tilde{\chi}(\Delta)=\sum_{j=-1}^{d-1}(-1)^j\beta_j(\Delta),
\]
where $d-1$ is the dimension of $\Delta$. We will also denote by $\chi(\Delta):=\tilde{\chi}(\Delta)+1$ the (unreduced) Euler characteristic of $\Delta$. If $\Delta$ is an orientable homology manifold without boundary, then Klee's formula
\begin{equation}\label{KleesFormula}
h_{d-i}(\Delta)-h_i(\Delta)=(-1)^i\binom{d}{i}((-1)^{d-1}\tilde{\chi}(\Delta)-1)
\end{equation}
(found in~\cite{Klee}) provides a relation between the symmetric entries in the $h$-vector of $\Delta$ in terms of $\tilde{\chi}(\Delta)$, abstracting the Dehn--Sommerville relations for simplicial polytopes. We will first examine the symmetric values of the dimensions of $\dim_\field \left(\field[\Delta]/\Sigma(\Theta;\field[\Delta])\right)_i^0$, in the case that $\Theta\subset A_1^0$.

\begin{theo}\label{DSTheorem}
Let $\Delta$ be a $(d-1)$-dimensional triangulation of a connected orientable manifold without boundary admitting a free group action by $G$. If $\Theta\subset A_1^0$ and $\widetilde{H}^{d-1}(\Delta)=\widetilde{H}^{d-1}(\Delta)^0$, then
\[
\dim_\field \left(\field[\Delta]/\Sigma(\Theta;\field[\Delta])\right)_i^0=\dim_\field \left(\field[\Delta]/\Sigma(\Theta;\field[\Delta])\right)_{d-i}^0.
\]
\end{theo}

\begin{proof}
If $\Delta=\{\emptyset\}$, then the result is immediate. Assume that $\Delta\not=\{\emptyset\}$. Applying Theorem~\ref{SigmaModuleThm}, we can write
\[
\dim_\field \left(\field[\Delta]/\Sigma(\Theta;\field[\Delta])\right)_i^0=\frac{h_i(\Delta)}{p}+\frac{(-1)^i(p-1)}{p}\binom{d}{i}+\binom{d}{i}\sum_{j=0}^i(-1)^{i-j-1}\beta_{j-1}(\Delta)^0.
\]
Using~\eqref{KleesFormula},
\begin{align*}
\frac{h_i(\Delta)}{p}-\frac{h_{d-i}(\Delta)}{p}&=\frac{(-1)^{i+1}}{p}\binom{d}{i}((-1)^{d-1}\tilde{\chi}(\Delta)-1)\\
&=\frac{(-1)^{d-i}}{p}\binom{d}{i}\tilde{\chi}(\Delta)+\frac{(-1)^i}{p}\binom{d}{i}\\
&=\frac{(-1)^{d-i}}{p}\binom{d}{i}\chi(\Delta)+\frac{(-1)^{d-i-1}}{p}\binom{d}{i}+\frac{(-1)^i}{p}\binom{d}{i}.
\end{align*}
Hence,
\begin{multline}\label{hVectorDifference}
\left[\frac{h_i(\Delta)}{p}+\frac{(-1)^i(p-1)}{p}\binom{d}{i}\right]-\left[\frac{h_{d-i}(\Delta)}{p}+\frac{(-1)^{d-i}(p-1)}{p}\binom{d}{i}\right]\\
=\frac{(-1)^{d-i}}{p}\binom{d}{i}\chi(\Delta)+(-1)^i\binom{d}{i}+(-1)^{d-i-1}\binom{d}{i}.
\end{multline}

Identify $\Delta$ with its geometric realization $|\Delta|$, and set $\Gamma= |\Delta|/G$. Since $G$ acts freely on $\Delta$, the projection $\pi:|\Delta|\to \Gamma$ is a covering space map. By~\cite[Proposition~3G.1]{Hatcher},
\[
H^i(|\Delta|)^0=H^i(|\Delta|)^G\cong H^i(\Gamma).
\]
Since $\widetilde{H}^{d-1}(\Delta)=\widetilde{H}^{d-1}(\Delta)^0$ by assumption, $\Gamma$ is itself an orientable manifold. Hence, Poincar\'e duality applies and $\beta_{j-1}(\Delta)^0=\beta_{d-j}(\Delta)^0$ for $j>1$ while $\beta_0(\Delta)^0=\beta_{d-1}(\Delta)^0-1$. Furthermore, $\beta_{-1}(\Delta)^0=0$ since $\Delta\not=\{\emptyset\}$. Then
\begin{align*}
\binom{d}{i}\sum_{j=0}^i(-1)^{i-j-1}\beta_{j-1}(\Delta)^0&=(-1)^{i+1}\binom{d}{i}+\binom{d}{i}\sum_{j=1}^i(-1)^{i-j-1}\beta_{d-j}(\Delta)^0\\
&=(-1)^{i+1}\binom{d}{i}+\binom{d}{i}\sum_{\ell=i+1}^{d}(-1)^{d-i-\ell}\beta_{\ell -1}(\Delta)^0.
\end{align*}
Thus,
\begin{multline*}
\binom{d}{i}\sum_{j=0}^i(-1)^{i-j-1}\beta_{j-1}(\Delta)^0-\binom{d}{i}\sum_{j=0}^{d-i}(-1)^{d-i-j-1}\beta_{j-1}(\Delta)^0\\
\begin{aligned}
&=(-1)^{i+1}\binom{d}{i}+\binom{d}{i}\sum_{j=0}^d(-1)^{d-i-j}\beta_{j-1}(\Delta)^0\\
&=(-1)^{i+1}\binom{d}{i}+(-1)^{d-i-1}\binom{d}{i}\tilde{\chi}(\Gamma).
\end{aligned}
\end{multline*}
Now note that $\Delta$ is a $p$-sheeted covering space for $\Gamma$. This implies that $\chi(\Delta)=p\chi(\Gamma)$ (see~\cite[Section~2.2]{Hatcher}). Hence, we can re-write the difference as
\begin{multline*}
(-1)^{i+1}\binom{d}{i}+(-1)^{d-i-1}\binom{d}{i}\tilde{\chi}(\Gamma)\\
\begin{aligned}
&=(-1)^{i+1}\binom{d}{i}+(-1)^{d-i-1}\binom{d}{i}\chi(\Gamma)+(-1)^{d-i}\binom{d}{i}\\
&=(-1)^{i+1}\binom{d}{i}+\frac{(-1)^{d-i-1}}{p}\binom{d}{i}\chi(\Delta)+(-1)^{d-i}\binom{d}{i}.
\end{aligned}
\end{multline*}
Adding this difference to~\eqref{hVectorDifference} completes the proof.
\end{proof}

\begin{rema}
In the case that $G=\ZZ/p\ZZ$ with $p$ odd, the hypothesis $\widetilde{H}^{d-1}(\Delta)=\widetilde{H}^{d-1}(\Delta)^0$ of Theorem~\ref{DSTheorem} is always true. Indeed, the fundamental class of $\Delta$ generating $\widetilde{H}^{d-1}(\Delta)$ is of the form
\[
z = \sum_{\widehat{\sigma}\in C^{d-1}(\Delta)}a_\sigma \widehat{\sigma}
\]
with $a_\sigma=\pm 1$ for all $\sigma$. If $h\in G$ acts on this class, then it can only flip the sign of some coefficients; it cannot introduce a factor of $\zeta^i$ across all coefficients for some non-trivial $i$.
\end{rema}

In fact, the theorem above can be greatly strengthened by appealing to algebraic properties of $\field[\Delta]/\Sigma(\Theta;\field[\Delta])$. First, we will need the following theorem (see~\cite[Theorem~1.4]{NS-Gorenstein} or~\cite[Remark~3.8]{BBMD}).

\begin{theo}
Let $\Delta$ be a triangulation of a $(d-1)$-dimensional connected manifold without boundary that is orientable over $\field$, and let $\Theta$ be a linear system of parameters for $\field[\Delta]$. Then $\field[\Delta]/\Sigma(\Theta;\field[\Delta])$ is an Artinian Gorenstein $\field$-algebra.
\end{theo}

With this fact, our next theorem quickly follows.

\begin{theo}\label{DSTheorem2}
Let $\Delta$ be a triangulation of a $(d-1)$-dimensional connected manifold without boundary that is orientable over $\field$ and admits a free group action by $G$, and let $\Theta$ be a $(\ZZ\times G)$-homogeneous linear system of parameters for $\field[\Delta]$. If $s$ is such that $\field[\Delta]/\Sigma(\Theta;\field[\Delta])_d$ is concentrated in $(\ZZ\times G)$-degree $(d, s)$, then
\[
\dim_\field \left(\field[\Delta]/\Sigma(\Theta;\field[\Delta])\right)_i^j = \dim_\field \left(\field[\Delta]/\Sigma(\Theta;\field[\Delta])\right)_{d-i}^{s-j}
\]
for $i=0, \ldots, d$ and any $j$, where we interpret $s-j$ modulo $p$.
\end{theo}

Before starting the proof of this theorem, note that $\dim_{\field}\field[\Delta]/\Sigma(\Theta;\field[\Delta])_d=1$ because $\field[\Delta]/\Sigma(\Theta;\field[\Delta])$ is Gorenstein and hence such an $s$ always exists.

\begin{proof}
Since $\field[\Delta]/\Sigma(\Theta;\field[\Delta])$ is Gorenstein, the product map
\[
\field[\Delta]/\Sigma(\Theta;\field[\Delta])_i \times \field[\Delta]/\Sigma(\Theta;\field[\Delta])_{d-i} \to \field[\Delta]/\Sigma(\Theta;\field[\Delta])_d
\]
is a perfect pairing (see~\cite[Theorem~2.79]{LefschetzProperties}). Furthermore, the finely-graded product maps
\[
\field[\Delta]/\Sigma(\Theta;\field[\Delta])_i^j \times \field[\Delta]/\Sigma(\Theta;\field[\Delta])_{d-i}^k \to \field[\Delta]/\Sigma(\Theta;\field[\Delta])_d^{j+k}
\]
only land in a non-zero component of $\field[\Delta]/\Sigma(\Theta;\field[\Delta])_d$ when $j+k\equiv s \pmod p$. That is,
\[
\field[\Delta]/\Sigma(\Theta;\field[\Delta])_i^j \times \field[\Delta]/\Sigma(\Theta;\field[\Delta])_{d-i}^{s-j} \to \field[\Delta]/\Sigma(\Theta;\field[\Delta])_d^{s}
\]
is a perfect pairing for each $i, j$ and hence
\[
\field[\Delta]/\Sigma(\Theta;\field[\Delta])_i^j \cong \field[\Delta]/\Sigma(\Theta;\field[\Delta])_{d-i}^{s-j}
\]
as vector spaces over $\field$.
\end{proof}

By comparing~\ref{DSTheorem} with~\ref{DSTheorem2}, we obtain the following corollary.

\begin{coro}
Let $\Delta$ be a $(d-1)$-dimensional triangulation of a connected orientable manifold without boundary admitting a free group action by $G$. If $\Theta\subset A_1^0$ and $\widetilde{H}^{d-1}(\Delta)=\widetilde{H}^{d-1}(\Delta)^0$, then $\field[\Delta]/\Sigma(\Theta;\field[\Delta])_d$ in concentrated in $(\ZZ\times G)$-degree $(d, 0)$ and
\[
\dim_\field \left(\field[\Delta]/\Sigma(\Theta;\field[\Delta])\right)_i^j = \dim_\field \left(\field[\Delta]/\Sigma(\Theta;\field[\Delta])\right)_{d-i}^{p-j}
\]
for $i=0, \ldots, d$ and $j=0, \ldots, p-1$.
\end{coro}


\section{Comments}\label{comments}

It was shown by Duval in~\cite{Duval} that a version of Theorem~\ref{Reisner} also holds when considering the face ring of a simplicial poset. Our main result then begs the following question.

\begin{question}
Does Theorem~\ref{ReisnerAnalog} extend to simplicial posets admitting free group actions?
\end{question}

An answer to this question will likely require a careful examination of a $(\ZZ\times G)$-grading imposed on the face ring of a simplicial poset, an object that at times is considerably more complicated than the Stanley--Reisner ring.

Stanley was able to examine some implications of equality being attained in his inequalities of Theorem~\ref{StanleyCMVeryFree} for the $G=\ZZ/2\ZZ$ case (see~\cite[Proposition~III.8.2]{St-96}). In light of these results, our next question naturally arises.

\begin{question}
If equality is attained in either inequality~\eqref{zeropart} or~\eqref{nonzeropart}, what can be said about the other entries in the $h$-vector of $\Delta$?
\end{question}

Lastly, this paper would feel incomplete without some mention of the $g$-conjecture. To this end, let $\Delta$ be a $(d-1)$-dimensional orientable $\field$-homology manifold and denote
\[
h_i''(\Delta)=\dim_\field \left(\field[\Delta]/\Sigma(\Theta;\field[\Delta])\right)_i=h_i(\Delta)+\binom{d}{i}\sum_{j=0}^i(-1)^{i-j-1}\beta_{j-1}(\Delta).
\]
Kalai's manifold $g$-conjecture asserts in part that $h_0''(\Delta)\le h_1''(\Delta)\le \cdots \le h_{\lfloor d/2\rfloor}''(\Delta)$ (in the case that $\Delta$ is a sphere, this is the same as the corresponding part of the usual $g$-conjecture). There is much evidence in favor of this statement, and it has been shown to be true in some special cases (see~\cite[Theorem~5.3]{Tight}, \cite[Theorem~1.5]{FaceNumbers}, and~\cite[Theorem~1.6]{NS-Gorenstein}). The inequality is sometimes demonstrated by exhibiting a Lefschetz element $\omega\in A_1$ such that the multiplication map
\[
\cdot\omega:\field[\Delta]/\Sigma(\Theta;\field[\Delta])_{i-1}\to\field[\Delta]/\Sigma(\Theta;\field[\Delta])_{i}
\]
is an injection for $1\le i\le \lfloor d/2\rfloor$. As has been our theme, we would like to show that similar inequalities hold under the finer grading. Indeed, if such an $\omega$ can in fact be chosen to be an element of $A_1^m$ for some $m$, then we immediately have the inequalities
\[
\dim_\field \left(\field[\Delta]/\Sigma(\Theta;\field[\Delta])\right)_{i-1}^j\le \dim_\field \left(\field[\Delta]/\Sigma(\Theta;\field[\Delta])\right)_i^{j+m}
\]
for any choice of $j$. Although the Lefschetz elements that have been found thus far cannot be specified to reside in some fixed homogeneous $(\ZZ\times G)$-degree of $A_1$, the strength of the finer grading on $\field[\Delta]/\Sigma(\Theta;\field[\Delta])$ prompts the following conjecture.

\begin{conjecture}
Let $\Delta$ be an orientable $\field$-homology manifold admitting a free group action by $\ZZ/p\ZZ$. Then there exists $m$ such that
\[
\dim_\field \left(\field[\Delta]/\Sigma(\Theta;\field[\Delta])\right)_{i-1}^j\le \dim_\field \left(\field[\Delta]/\Sigma(\Theta;\field[\Delta])\right)_i^{j+m}
\]
for $1\le i\le \lfloor d/2\rfloor$ and $0\le j\le p-1$.
\end{conjecture}

\longthanks The author would like to thank Martina Juhnke-Kubitzke, Isabella Novik, Travis Scholl, and two anonymous referees for many helpful comments and conversations.

\begin{thebibliography}{25}
\bibitem[1]{AdinThesis} Ron Adin, Combinatorial structure of simplicial complexes with symmetry, Ph.D. thesis, Hebrew University, Jerusalem, 1991.
\bibitem[2]{Bourbaki} Nicolas Bourbaki, Algebra. I. Chapters 1–3, Elements of Mathematics, Springer, 1989, Translated from the French, Reprint of the 1974 edition.
\bibitem[3]{Duval} Art M. Duval, Free resolutions of simplicial posets, J. Algebra 188 (1997), no. 1, 363–399.
\bibitem[4]{Garsia} Adriano M. Garsia and Dennis Stanton, Group actions of Stanley-Reisner rings and invariants of permutation groups, Adv. Math. 51 (1984), no. 2, 107–201.
\bibitem[5]{Goto} Shiro Goto, On the associated graded rings of parameter ideals in Buchsbaum rings, J. Algebra 85 (1983), no. 2, 490–534.
\bibitem[6]{Grabe} Hans-Gert Gräbe, The canonical module of a Stanley-Reisner ring, J. Algebra 86 (1984), no. 1, 272–281.
\bibitem[7]{GrabeDS} Hans-Gert Gräbe, Generalized Dehn-Sommerville equations and an upper bound theorem, Beitr. Algebra Geom. (1987), no. 25, 47–60.
\bibitem[8]{LefschetzProperties} Tadahito Harima, Toshiaki Maeno, Hideaki Morita, Yasuhide Numata, Akihito Wachi, and Junzo Watanabe, The Lefschetz properties, Lecture Notes in Math., vol. 2080, Springer, 2013.
\bibitem[9]{Hatcher} Allen Hatcher, Algebraic topology, Cambridge University Press, 2002.
\bibitem[10]{24Hours} Srikanth B. Iyengar, Graham J. Leuschke, Anton Leykin, Claudia Miller, Ezra Miller, Anurag K. Singh, and Uli Walther, Twenty-four hours of local cohomology, Graduate Studies in Mathematics, vol. 87, American Mathematical Society, 2007.
\bibitem[11]{Klee} Victor Klee, A combinatorial analogue of Poincaré’s duality theorem, Canad. J. Math. 16 (1964), 517–531.
\bibitem[12]{characterizations} Mitsuhiro Miyazaki, Characterizations of Buchsbaum complexes, Manuscr. Math. 63 (1989), no. 2, 245–254.
\bibitem[13]{canonicalMap} Mitsuhiro Miyazaki, On the canonical map to the local cohomology of a Stanley-Reisner ring, Bull. Kyoto Univ. Educ., Ser. B (1991), no. 79, 1–8.
\bibitem[14]{Tight} Satoshi Murai, Tight combinatorial manifolds and graded Betti numbers, Collect. Math. 66 (2015), no. 3, 367–386.
\bibitem[15]{FaceNumbers} Satoshi Murai and Isabella Novik, Face numbers of manifolds with boundary, Int. Math. Res. Not. (2017), no. 12, 3603–3646.
\bibitem[16]{BBMD} Satoshi Murai, Isabella Novik, and Ken-ichi Yoshida, A duality in Buchsbaum rings and triangulated manifolds, Algebra Number Theory 11 (2017), no. 3, 635–656.
\bibitem[17]{NS-Gorenstein} Isabella Novik and Ed Swartz, Gorenstein rings through face rings of manifolds, Compos. Math. 145 (2009), no. 4, 993–1000.
\bibitem[18]{NS-Socle} Isabella Novik and Ed Swartz, Socles of Buchsbaum modules, complexes and posets, Adv. Math. 222 (2009), no. 6, 2059–2084.
\bibitem[19]{Reiner} Victor Reiner, Quotients of Coxeter complexes and P -partitions, Mem. Am. Math. Soc. 95 (1992), no. 460, vi+134.
\bibitem[20]{Reisner} Gerald Allen Reisner, Cohen-Macaulay quotients of polynomial rings, Adv. Math. 21 (1976), no. 1, 30–49.
\bibitem[21]{Schenzel} Peter Schenzel, On the number of faces of simplicial complexes and the purity of Frobenius, Math. Z. 178 (1981), no. 1, 125–142.
\bibitem[22]{St-87} Richard P. Stanley, On the number of faces of centrally-symmetric simplicial polytopes, Graphs Combin. 3 (1987), no. 1, 55–66.
\bibitem[23]{StanleyPosets} Richard P. Stanley, f -vectors and h-vectors of simplicial posets, J. Pure Appl. Algebra 71 (1991), no. 2-3, 319–331.
\bibitem[24]{St-96} Richard P. Stanley, Combinatorics and commutative algebra, second ed., Progress in Mathematics, vol. 41, Birkhäuser, 1996.
\bibitem[25]{StVo} Jürgen Stückrad and Wolfgang Vogel, Buchsbaum rings and applications. An interaction between algebra, geometry, and topology, Mathematische Monographien, vol. 21, VEB Deutscher Verlag der Wissenschaften, 1986.
\end{thebibliography}
\end{document}
